\documentclass[11pt,fontset=fandol]{ctexart}
\usepackage[a4paper,margin=23mm]{geometry}
\usepackage{amsmath,amssymb,mathtools,bm,mathrsfs,esint}
\usepackage{graphicx,longtable,booktabs,array,calc}
\usepackage[normalem]{ulem}
\usepackage{xcolor,hyperref}
\hypersetup{hidelinks}
\setcounter{secnumdepth}{-1}
\newcounter{none}
\providecommand{\tightlist}{\setlength{\itemsep}{0pt}\setlength{\parskip}{0pt}}
\providecommand{\passthrough}[1]{#1}
\usepackage{amssymb,mathrsfs,esint,bm,newunicodechar}
\allowdisplaybreaks
\setlength{\emergencystretch}{3em}
\providecommand{\llangle}{\langle\!\langle}
\providecommand{\rrangle}{\rangle\!\rangle}
\newunicodechar{∞}{\ensuremath{\infty}}
\newunicodechar{∈}{\ensuremath{\in}}
\newunicodechar{∑}{\ensuremath{\sum}}
\newunicodechar{∫}{\ensuremath{\int}}
\newunicodechar{≤}{\ensuremath{\leq}}
\newunicodechar{≥}{\ensuremath{\geq}}
\newunicodechar{≠}{\ensuremath{\neq}}
\newunicodechar{→}{\ensuremath{\to}}
\newunicodechar{⇒}{\ensuremath{\Rightarrow}}
\newunicodechar{α}{\ensuremath{\alpha}}
\newunicodechar{β}{\ensuremath{\beta}}
\newunicodechar{γ}{\ensuremath{\gamma}}
\newunicodechar{λ}{\ensuremath{\lambda}}
\newunicodechar{φ}{\ensuremath{\varphi}}
\newunicodechar{ε}{\ensuremath{\varepsilon}}
\newunicodechar{σ}{\ensuremath{\sigma}}
\newunicodechar{θ}{\ensuremath{\theta}}
\newunicodechar{δ}{\ensuremath{\delta}}
\newunicodechar{η}{\ensuremath{\eta}}
\newunicodechar{π}{\ensuremath{\pi}}
\newunicodechar{ξ}{\ensuremath{\xi}}
\newunicodechar{Σ}{\ensuremath{\Sigma}}
\newunicodechar{Δ}{\ensuremath{\Delta}}
\newunicodechar{∂}{\ensuremath{\partial}}
\newunicodechar{∀}{\ensuremath{\forall}}
\newunicodechar{∃}{\ensuremath{\exists}}
\newunicodechar{√}{\ensuremath{\surd}}
\newunicodechar{−}{\ensuremath{-}}
\newunicodechar{μ}{\ensuremath{\mu}}
\newunicodechar{ρ}{\ensuremath{\rho}}
\newunicodechar{τ}{\ensuremath{\tau}}
\newunicodechar{ω}{\ensuremath{\omega}}
\newunicodechar{Ω}{\ensuremath{\Omega}}
\newunicodechar{χ}{\ensuremath{\chi}}
\newunicodechar{ψ}{\ensuremath{\psi}}
\newunicodechar{Γ}{\ensuremath{\Gamma}}
\newunicodechar{Φ}{\ensuremath{\Phi}}
\newunicodechar{Λ}{\ensuremath{\Lambda}}
\newunicodechar{ℝ}{\ensuremath{\mathbb{R}}}
\newunicodechar{ℤ}{\ensuremath{\mathbb{Z}}}
\newunicodechar{ℕ}{\ensuremath{\mathbb{N}}}
\newunicodechar{△}{\ensuremath{\triangle}}
\newunicodechar{⊗}{\ensuremath{\otimes}}
\begin{document}
\section{2022年高等代数真题}\label{ux5e74ux9ad8ux7b49ux4ee3ux6570ux771fux9898}

来源：\texttt{HDU/601数学分析/2022/2022年杭州电子科技大学数分高代真题及答案(双科合卷).pdf}。原
PDF 页序：1, 2。

\section{杭州电子科技大学2022年高等代数真题}\label{ux676dux5ddeux7535ux5b50ux79d1ux6280ux5927ux5b662022ux5e74ux9ad8ux7b49ux4ee3ux6570ux771fux9898}

\begin{quote}
整理说明：本文保留原卷全部七题和解析的数学内容，清理混入公式的水印并统一矩阵语法。原解析第一题的中间行列式被水印严重破坏，以下根据清楚的题面重建推导并明确标出；第六、七题缺损矩阵的数值依据题面、后续答案和直接运算复核恢复。不能将OCR产生的额外行列理解为原矩阵的组成部分。
\end{quote}

\subsection{一、计算行列式（15分）}\label{ux4e00ux8ba1ux7b97ux884cux5217ux5f0f15ux5206}

\[
D_n=\begin{vmatrix}
x&x&\cdots&x&x&a_1\\
x&x&\cdots&x&a_2&x\\
x&x&\cdots&a_3&x&x\\
\vdots&\vdots&&\vdots&\vdots&\vdots\\
x&a_{n-1}&\cdots&x&x&x\\
a_n&x&\cdots&x&x&x
\end{vmatrix}.
\]

\subsection{二、多项式互素（15分）}\label{ux4e8cux591aux9879ux5f0fux4e92ux7d2015ux5206}

设 \(f(x),g(x)\) 为数域 \(P\) 上的多项式，\(a,b,c,d\in P\)，且
\(ad-bc\ne0\)。证明 \[
(af(x)+bg(x),cf(x)+dg(x))=1
\quad\Longleftrightarrow\quad(f(x),g(x))=1.
\]

\subsection{三、齐次线性方程组（20分）}\label{ux4e09ux9f50ux6b21ux7ebfux6027ux65b9ux7a0bux7ec420ux5206}

\[
\begin{cases}
(a_1+b)x_1+a_2x_2+a_3x_3+a_4x_4=0,\\
a_1x_1+(a_2+b)x_2+a_3x_3+a_4x_4=0,\\
a_1x_1+a_2x_2+(a_3+b)x_3+a_4x_4=0,\\
a_1x_1+a_2x_2+a_3x_3+(a_4+b)x_4=0,
\end{cases}
\qquad a_1+a_2+a_3+a_4\ne0.
\]

（1）参数满足何种关系时方程组仅有零解？

（2）参数满足何种关系时方程组有非零解？此时求一个基础解系。

\subsection{四、幂等矩阵（每小题10分，共30分）}\label{ux56dbux5e42ux7b49ux77e9ux9635ux6bcfux5c0fux989810ux5206ux517130ux5206}

设 \(A,B\) 是数域 \(P\) 上的 \(n\) 阶矩阵，\(A^2=A,B^2=B\)。证明：

（1）\(r(A)+r(A-E_n)=n\)。

（2）\(A\) 相似于 \(\begin{pmatrix}E_r&0\\0&0\end{pmatrix}\)，其中
\(r=r(A)\)。

（3）若 \(AB=BA\)，存在可逆矩阵 \(P\) 使 \(P^{-1}AP\) 与 \(P^{-1}BP\)
同时为对角矩阵。

\subsection{五、相似与合同（每小题10分，共20分）}\label{ux4e94ux76f8ux4f3cux4e0eux5408ux540cux6bcfux5c0fux989810ux5206ux517120ux5206}

\begin{enumerate}
\def\labelenumi{\arabic{enumi}.}
\item
  设 \(A,B\) 为 \(n\) 阶方阵，证明相似矩阵有相同特征值；反之是否成立？
\item
  设 \[
  A=\begin{pmatrix}2&-1&-1\\-1&2&-1\\-1&-1&2\end{pmatrix},\qquad
  B=\begin{pmatrix}2&1&0\\1&2&0\\0&0&0\end{pmatrix}.
  \] 判断二者是否相似、是否合同，并说明理由。
\end{enumerate}

\subsection{六、线性变换的矩阵、核与值域（每小题10分，共20分）}\label{ux516dux7ebfux6027ux53d8ux6362ux7684ux77e9ux9635ux6838ux4e0eux503cux57dfux6bcfux5c0fux989810ux5206ux517120ux5206}

设 \(\varepsilon_1,\varepsilon_2,\varepsilon_3,\varepsilon_4\)
是四维线性空间 \(V\) 的一组基，线性变换 \(\sigma\) 在该基下的矩阵为 \[
A=\begin{pmatrix}1&2&3&-6\\3&2&1&-6\\2&3&1&6\\1&1&1&-3\end{pmatrix}.
\]

（1）求 \(\sigma\) 在基 \[
\eta_1=\varepsilon_1+\varepsilon_2,\quad
\eta_2=\varepsilon_2-\varepsilon_3,\quad
\eta_3=\varepsilon_3,\quad\eta_4=\varepsilon_4
\] 下的矩阵。

（2）求 \(\sigma\) 的核与值域各一组基。

\subsection{七、二次型（每小题10分，共30分）}\label{ux4e03ux4e8cux6b21ux578bux6bcfux5c0fux989810ux5206ux517130ux5206}

设 \(f(x_1,x_2,x_3)=x^TAx\)，其中 \[
x=\begin{pmatrix}x_1\\x_2\\x_3\end{pmatrix},\qquad
A=\begin{pmatrix}b&-2&0\\-2&a&-2\\0&-2&0\end{pmatrix}.
\]

（1）当 \(A\) 有特征值 \(1,-2,4\) 时，求 \(a,b\)。

（2）对求得的 \(a,b\)，用正交线性替换 \(X=QY\)
将二次型化为标准形，并求正交矩阵 \(Q\)。

（3）问 \(t\) 为何值时 \(A+tE\) 正定、半正定？
\end{document}
